\section{Configuraciones}

\subsection{Requisitos Previos}

\begin{itemize}[leftmargin=*]
    \item Docker Desktop corriendo
    \item Node.js 18+ instalado localmente
    \item Haber ejecutado \texttt{npm install} dentro del directorio \texttt{perf/}
\end{itemize}

\subsection{Levantar el stack}

\begin{verbatim}
# Desde la raíz del proyecto
docker compose up --build -d
\end{verbatim}

Verificar que todos los contenedores estén healthy:

\begin{verbatim}
docker compose ps
\end{verbatim}

Todos deben aparecer en estado \texttt{healthy} o \texttt{running}. Si \texttt{app1/app2/app3} aparece como \texttt{starting}, esperar \(\sim\)15 segundos y volver a verificar.

Verificar que la API responde:

\begin{verbatim}
curl http://localhost:5555/health
# Esperado: {"status":"ok"}
\end{verbatim}

\subsection{Configurar Grafana (primera vez)}

\begin{enumerate}[leftmargin=*]
    \item Abrir \texttt{http://localhost:80} en el browser (usuario/contraseña: \texttt{admin/admin}).
    \item Ir a \textbf{Connections → Data Sources → Add data source}.
    \item Elegir \textbf{Graphite}.
    \item En el campo URL ingresar: \texttt{http://graphite:80}
    \item Hacer click en \textbf{Save \& test} — debe aparecer ``Data source is working''.
    \item Ir a \textbf{Dashboards → Import}.
    \item Hacer click en \textbf{Upload dashboard JSON file} y seleccionar \texttt{perf/dashboard.json}.
    \item Asignar el datasource Graphite y hacer \textbf{Import}.
\end{enumerate}

El dashboard queda guardado permanentemente en el volumen de Grafana.

\subsection{Limpiar y reiniciar el stack entre pruebas}

Si los saldos de las cuentas se agotan durante las pruebas de stress (el exchange empieza a devolver 500), se puede reiniciar el estado:

\begin{verbatim}
# Reiniciar solo la app (Redis mantiene el estado si se configuró persistencia)
docker compose restart app1 app2 app3
\end{verbatim}

Para empezar desde cero:

\begin{verbatim}
docker compose down -v   # borra volúmenes incluyendo Redis
docker compose up --build -d
\end{verbatim}
